\section{食品安全法}
\begin{enumerate}
  \item 比重越来越重。
  \item 范围。生产经营、贮存和运输、添加剂、相关产品的安全管理（包装）、\zd{初级食用农产品}（大米、大豆）\yjt 农产品质量安全法。
  \item 食品安全风险监测。中央和省级。
  \item 风险评估制度。采样；不得收取费用。\sjt 简单了解。
  \item \zd{食品安全标准}。强制执行的标准。
        \begin{enumerate}
          \item 食品添加剂。品种、使用范围、用量。
          \item 专供婴幼儿和其他特定人群的主辅食品的营养成分要求。
          \item 标签、标志、说明书。
          \item 3 类标准。
                \begin{enumerate}
                  \item 国家标准。两部门。国务院卫生行政部门会同国务院食品安全监管部。
                  \item 地方标准。地方特色食品。\al{云南鲜花饼}。\zd{不是全国通用 + 国标缺失}。有国标后 \yjt \zd{不能并存，即行废止}。
                  \item 企业标准。\zd{严于}国标 \hz 地方；省级备案。
                \end{enumerate}
        \end{enumerate}
  \item 食品生产经营的一般规定。营业执照 + 食品生产经营许可证。
        \begin{enumerate}
          \item 销售食用农产品不要求取得许可的。
          \item \zd{对集中交易市场的管理}（集贸市场、农贸市场）。审查许可证 + 检查 \yjt 及时制止，立即报告。
          \item 网络食品交易第三方平台提供者。实名登记 + 审查其许可证 \yjt 及时制止，立即报告。
        \end{enumerate}
  \item 召回制度。
        \begin{enumerate}
          \item 生产者召回是原则，经营者召回是例外（经营者原因导致）。该召回但没有召回的 \yjt 强制召回。
          \item 无害化处理 \hz 销毁。
          \item 标签、标志、说明书不符合安全标准而被召回是可以继续销售（产品本身没问题）。
        \end{enumerate}
  \item 散装、预包装（预先 + 定量）食品也要打标签；禁止相关行政部门推荐食品。
  \item 几类特殊食品。
        \begin{enumerate}
          \item \zd{保健食品}。
                \begin{enumerate}
                  \item 保健食品原料目录。
                        \begin{enumerate}
                          \item 3 部门。国务院食品药品监督管理部门会同卫生行政部门、中医药管理部门制定。
                          \item 进口保健食品原料要注册；其余视情况备案。
                        \end{enumerate}
                  \item 保健功能目录。
                  \item 标签、说明书\zd{不得涉及疾病预防、治疗}功能。广告当中还应当声明\zd{本品不能替代药品}。
                \end{enumerate}
          \item 婴幼儿配方食品。
          \item 特殊医学用途配方食品。国务院食品安全监督管理部门注册，省备案 + 公布目录。
        \end{enumerate}
  \item 食品的检验。抽样检验；县级以上食品安全监督管理部门；不能收费；\zd{不得免检}。
  \item 食品安全事故的处置。
        \begin{enumerate}
          \item 应急预案。不做考试要求。
                \begin{enumerate}
                  \item 国务院组织制定国家食品安全事故的应急预案；地方\zd{政府}制定。
                  \item 处置方案企业制定。
                \end{enumerate}
          \item 双报告制度。事故单位 + 医院 \yjt 卫生行政部门、县食品安全监督管理部门 \yjt 本级、上级政府。
          \item 流行病学调查（流调）。疾控中心。双报告。
        \end{enumerate}
  \item 民事责任。不符合食品安全标准。找市场监督管理局。\zd{消费者证明损害和初步因果关系}。
        \begin{enumerate}
          \item 首负责任制。找到谁，谁先先行赔付，不得推诿。
          \item 经营者。进货已查验 + 不知情可免予行政处罚。
          \item 虚假广告、网络平台连带责任。
          \item 标签、说明书有瑕疵（\al{错别字}），但不影响食品安全 + 不会造成误导除外。
          \item \zd{食品、药品支持知假买假}；赠品不得以未支付对价抗辩。
          \item 不达标\zd{惩罚性赔偿}。价款 10 倍 \hz 损失 3 倍，1000 保底。生产者不要求明知，销售者要求明知。
        \end{enumerate}
\end{enumerate}
